%% kapa: exam=1 round=32 year=2021 period=1 form=A subject=civil_law subject_ko=민법 seq=15 number=15
%% source: Q-Net 감정평가사 기출문제 (cst00309) artlSeq=5208397 file=extracted/1008/5208397/A형/2021년도 제32회 감정평가사 1차 1교시 A형.pdf
%% conversion: kordoc 4.24.51 -> tools/gen1.py (auto); review=proofread-sonnet; flags=none
\begin{question}{15}
복대리에 관한 설명으로 옳은 것은? (다툼이 있으면 판례에 따름)
\begin{choices}
  \item 복대리인은 제3자에 대하여 대리인과 동일한 권리의무가 있다.
  \item 본인의 묵시적 승낙에 기초한 임의대리인의 복임권행사는 허용되지 않는다.
  \item 임의대리인이 본인의 명시적 승낙을 얻어 복대리인을 선임한 때에는 본인에 대하여 그 선임감독에 관한 책임이 없다.
  \item 법정대리인이 그 자신의 이름으로 선임한 복대리인은 법정대리인의 대리인이다.
  \item 복대리인의 대리행위에 대해서는 표현대리가 성립할 수 없다.
\end{choices}
\end{question}
\answer{1}
